% 来源：1008-9.pdf 全部；1013-10.pdf 第 1--2 页。
\section{密度矩阵方法}
\subsection{纯态、混合态与密度算符}
纯态可由一个归一化态矢量 $\Ket{\psi}$ 描述。若只知道系统以概率 $p_s$ 制备在归一化态 $\Ket{\psi_s}$ 中，则用统计系综
\begin{equation}
\left\{p_s,\Ket{\psi_s}\right\},\qquad p_s\geq0,\qquad\sum_s p_s=1
\end{equation}
描述。参与态不必正交，参与态数目也可超过 Hilbert 空间维数。相干叠加 $\sum_s c_s\Ket{\psi_s}$ 保留振幅的相对相位，统计混合则对各次制备的测量结果作概率平均。定义密度算符
\begin{equation}
\boxed{\hat\rho=\sum_s p_s\Ket{\psi_s}\Bra{\psi_s}.}
\label{eq:density-ensemble}
\end{equation}
纯态对应 $\hat\rho=\Ket{\psi}\Bra{\psi}$。有多个非零权重并不足以保证混合态：若参与态仅相差整体相位，仍得到同一个纯态投影。以下主要讨论有限维系统；无穷维情形要求密度算符为正的迹类算符，所用期望值存在。

在正交归一基底 $\left\{\Ket{a_n}\right\}$ 中插入两次完备性，得到
\begin{align}
\hat\rho&=\sum_{n,m}\Ket{a_n}\rho_{nm}\Bra{a_m},\\
\rho_{nm}&=\Braket{a_n|\hat\rho|a_m}
=\sum_s p_s\Braket{a_n|\psi_s}\Braket{\psi_s|a_m}.
\end{align}
$\hat\rho$ 是抽象算符，$\rho$ 是选定基底中的密度矩阵。若纯态的系数列向量为 $c$，则 $\rho=cc^\dagger$。非简并可观测量 $\hat A$ 的本征基中，$\rho_{nn}$ 是测得 $a_n$ 的概率；若简并，应对相应本征子空间内的对角元求和。非对角元描述相对于所选基底的相干性，不能直接当作概率。

\subsection{基本性质与纯度判据}
由定义取共轭、求迹，并对任意 $\Ket{\phi}$ 取矩阵元，可得
\begin{align}
\hat\rho^\dagger&=\hat\rho,\\
\mathrm{Tr}\,\hat\rho&=\sum_s p_s\Braket{\psi_s|\psi_s}=1,\\
\Braket{\phi|\hat\rho|\phi}&=\sum_s p_s\left|\Braket{\phi|\psi_s}\right|^2\geq0.
\end{align}
密度算符是迹为一的半正定厄米算符。反之，满足这些条件的算符可谱分解为
\begin{equation}
\hat\rho=\sum_\alpha\lambda_\alpha\Ket{r_\alpha}\Bra{r_\alpha},\qquad
\lambda_\alpha\geq0,\qquad\sum_\alpha\lambda_\alpha=1,
\end{equation}
因而能描述一个统计系综，且 $0\leq\lambda_\alpha\leq1$。定义纯度 $\mathcal P=\mathrm{Tr}\,\hat\rho^2$，在 $d$ 维空间中有
\begin{equation}
\boxed{\frac1d\leq\mathcal P=\sum_\alpha\lambda_\alpha^2\leq1.}
\end{equation}
上界来自 $\lambda_\alpha^2\leq\lambda_\alpha$；下界来自 Cauchy--Schwarz 不等式
$\left(\sum_\alpha\lambda_\alpha\right)^2\leq d\sum_\alpha\lambda_\alpha^2$。
上界取等号要求本征值均为零或一，由归一化知恰有一个本征值为一。因此
\begin{equation}
\boxed{\text{纯态}\ \Longleftrightarrow\ \hat\rho^2=\hat\rho
\ \Longleftrightarrow\ \mathrm{Tr}\,\hat\rho^2=1.}
\end{equation}
混合态满足 $\mathcal P<1$，其所有本征值均小于一；最大混合态为 $\hat I/d$。对任意系综分解也可直接计算
\begin{equation}
\mathcal P=\sum_{s,r}p_sp_r\left|\Braket{\psi_s|\psi_r}\right|^2.
\end{equation}
若非零权重的参与态中至少有两条不同射线，对应内积的模严格小于一，才可据此推出严格不等式。

\subsection{期望值与时间演化}
令 $\hat B$ 为可观测量，在系综平均中插入完备性，
\begin{align}
\overline B&=\sum_s p_s\Braket{\psi_s|\hat B|\psi_s}\notag\\
&=\sum_{s,n}p_s\Braket{a_n|\psi_s}\Braket{\psi_s|\hat B|a_n}
=\boxed{\mathrm{Tr}\left(\hat\rho\hat B\right).}
\end{align}
这个表达式与表象无关。两个制备系综只要给出相同的密度算符，就给出相同的全部单次测量统计。

对封闭系统，假设各参与态服从同一个 Schr\"odinger 方程且 $p_s$ 不随时间改变。由
\begin{equation}
\mathrm{i}\hbar\frac{\mathrm{d}}{\mathrm{d}t}\Ket{\psi_s}=\hat H\Ket{\psi_s},\qquad
\mathrm{i}\hbar\frac{\mathrm{d}}{\mathrm{d}t}\Bra{\psi_s}=-\Bra{\psi_s}\hat H
\end{equation}
和投影算符的乘积求导，得到 von Neumann 方程及其解
\begin{equation}
\boxed{\mathrm{i}\hbar\frac{\mathrm{d}\hat\rho}{\mathrm{d}t}=\left[\hat H,\hat\rho\right],}
\qquad
\hat\rho\left(t\right)=\hat U\left(t,t_0\right)\hat\rho\left(t_0\right)\hat U^\dagger\left(t,t_0\right).
\end{equation}
由幺正性及迹的循环性，
\begin{equation}
\mathrm{Tr}\,\hat\rho^2\left(t\right)
=\mathrm{Tr}\left[\hat U\hat\rho^2\left(t_0\right)\hat U^\dagger\right]
=\mathrm{Tr}\,\hat\rho^2\left(t_0\right).
\end{equation}
封闭系统的幺正演化保持密度算符的谱与纯度；子系统的约化态一般不满足这个封闭演化方程。

\subsection{自旋二分之一的例子}
取 $\hat S_z$ 的有序本征基 $\Ket{\uparrow},\Ket{\downarrow}$，自旋算符的矩阵表示为
\begin{equation}
S_x=\frac\hbar2\left(\begin{matrix}0&1\\1&0\end{matrix}\right),\quad
S_y=\frac\hbar2\left(\begin{matrix}0&-\mathrm{i}\\\mathrm{i}&0\end{matrix}\right),\quad
S_z=\frac\hbar2\left(\begin{matrix}1&0\\0&-1\end{matrix}\right).
\end{equation}
例如 $\hat S_x=\hbar\left(\Ket{\uparrow}\Bra{\downarrow}+\Ket{\downarrow}\Bra{\uparrow}\right)/2$。对纯态
\begin{equation}
\Ket{\psi_\phi}=\frac1{\sqrt2}\left(\mathrm{e}^{\mathrm{i}\phi}\Ket{\uparrow}+\Ket{\downarrow}\right),\qquad
\rho_\phi=\frac12\left(\begin{matrix}1&\mathrm{e}^{\mathrm{i}\phi}\\\mathrm{e}^{-\mathrm{i}\phi}&1\end{matrix}\right),
\end{equation}
矩阵乘法给出
\begin{equation}
\rho_\phi S_x=\frac\hbar4\left(\begin{matrix}\mathrm{e}^{\mathrm{i}\phi}&1\\1&\mathrm{e}^{-\mathrm{i}\phi}\end{matrix}\right),\qquad
\overline S_x=\frac\hbar2\cos\phi.
\end{equation}
另有 $\overline S_y=-\hbar\sin\phi/2$、$\overline S_z=0$。等概率制备上下自旋的混合态则为
\begin{equation}
\hat\rho_{\mathrm m}=\frac12\Ket{\uparrow}\Bra{\uparrow}+\frac12\Ket{\downarrow}\Bra{\downarrow}=\frac{\hat I}{2}.
\end{equation}
其三个自旋分量平均值均为零，纯度为 $1/2$。在 $\phi=0$ 时，两者测量 $S_z$ 的分布相同，测量 $S_x$ 却能区分它们，原因是纯态中的相干交叉项仍存在。

\subsection{系综分解不唯一与表象依赖}
记 $\Ket{\pm}_x=\left(\Ket{\uparrow}\pm\Ket{\downarrow}\right)/\sqrt2$，
$\Ket{\pm}_y=\left(\Ket{\uparrow}\pm\mathrm{i}\Ket{\downarrow}\right)/\sqrt2$。任意一组正交归一的二能级基底都给出 $\hat I/2$ 的等权分解，特别是 $x$、$y$、$z$ 三组自旋基底。
即使本征值不简并，非正交参与态的分解仍不唯一。例如
\begin{align}
\hat\rho&=\frac13\Ket{\uparrow}\Bra{\uparrow}+\frac23\Ket{\downarrow}\Bra{\downarrow}\notag\\
&=\frac12\Ket{\psi_+}\Bra{\psi_+}+\frac12\Ket{\psi_-}\Bra{\psi_-}\notag\\
&=\frac49\Ket{\varphi_1}\Bra{\varphi_1}+\frac59\Ket{\varphi_2}\Bra{\varphi_2},
\end{align}
其中
\begin{align}
\Ket{\psi_\pm}&=\frac{\Ket{\uparrow}\pm\sqrt2\Ket{\downarrow}}{\sqrt3},\\
\Ket{\varphi_1}&=\frac{\Ket{\uparrow}+\Ket{\downarrow}}{\sqrt2},\qquad
\Ket{\varphi_2}=\frac{\Ket{\uparrow}-2\Ket{\downarrow}}{\sqrt5}.
\end{align}
展开投影算符可见交叉项相消，三式相同。因此不能从 $\hat\rho$ 唯一反推出制备过程。
密度矩阵是否对角也取决于基底：上述例子在 $z$、$x$ 基底分别给出
\begin{equation}
\rho_z=\left(\begin{matrix}1/3&0\\0&2/3\end{matrix}\right),\qquad
\rho_x=\left(\begin{matrix}1/2&-1/6\\-1/6&1/2\end{matrix}\right).
\end{equation}
混合态并非在任意表象下都没有非对角元。另一个课堂例子是
\begin{equation}
\hat\rho=\frac12\Ket{\uparrow}\Bra{\uparrow}+\frac14\Ket{\downarrow}\Bra{\downarrow}+\frac14\Ket{+}_x\Bra{+}_x
\ \Longrightarrow\ 
\rho_z=\left(\begin{matrix}5/8&1/8\\1/8&3/8\end{matrix}\right).
\end{equation}
